Joint-embedding predictive architectures (JEPAs) learn world models by predicting future
representations rather than future pixels, and the objective alone decides what those
representations keep. We show that in stochastic worlds this freedom works against the controller
through two independent mechanisms. The first is representational: the squared next-embedding loss
rewards whatever is predictable, so when the controlled state carries fresh unobservable innovations
and a nuisance does not, the objective learns to ignore the state. In a matched three-seed study on
a minimal red-cube world, the same spatial JEPA that localizes the cube to $0.12$ pixels under
per-frame background noise loses it entirely ($21.3$ pixels) once hidden wind is added, while
persistent noise under the same wind keeps it at $0.67$ pixels: the effect of nuisance persistence
reverses with the dynamics. The second mechanism is predictive: a squared loss yields the
conditional mean, and two exact examples show that a correct mean can misrank actions and that a
correct one-step mean need not compose into a multi-step mean. In a gate world whose safe passage
is revealed late, the JEPA's frozen latent supports near-perfect survival when a policy is trained
on real transitions through its own decoder, yet policies trained inside its imagination
\pendingresult{gate imagination outcome, from the shared-CNN campaign}. In Craftax under matched
data and training budgets, \pendingresult{tool-acquisition outcome, from the inventory campaign}.
These results establish that avoiding representation collapse and accurately predicting latent
states are insufficient for successful stochastic control: distributional dynamics repair the
second mechanism only when the representation has survived the first.
